% =============================================================================
% Project : THE_PHILOSOPHY_OF_VIABILITY_LIBRARY_PROJECT
% Book    : Book I — Ontology of Viability
% Book ID : VS.B01
% Part    : Part I — Identity and Manifestation
% Chapter : Chapter 08 — Comparative Philosophical Examination
% Chapter ID : VS.B01.P01.C08
%
% File :
%   4_ASSETS/tables/VS.B01.P01.C08.S05.SS05.P005.TL01.aristotelian_matter_form_composite.tex
%
% Section / Chapter / Asset:
%   Aristotelian Matter, Form, and Composite
%
% Structural Identifier :
%   VS.B01.P01.C08.S05.SS05.P005.TL01
%
% Asset Type : TL
%
% Purpose :
%   To distinguish the explanatory roles of matter, form, and
%   composite in Aristotle's hylomorphic analysis.
%
% =============================================================================

\begin{table}[htbp]
\centering
\caption{Matter, Form, and Composite in Aristotle's Hylomorphic Analysis}
\label{tab:aristotle-matter-form-composite}

\small
\renewcommand{\arraystretch}{1.3}

\begin{tabularx}{\textwidth}{
    >{\raggedright\arraybackslash}p{0.17\textwidth}
    >{\raggedright\arraybackslash}p{0.37\textwidth}
    >{\raggedright\arraybackslash}X
}
\toprule
\textbf{Aspect} &
\textbf{Explanatory Role} &
\textbf{Interpretive Limit}
\\
\midrule

Matter (\textit{hylē})
&
That from which a substance is constituted, or the
underlying subject considered in relation to a process
of change.
&
Matter considered without substantial determination
does not by itself explain why a concrete substance
is the kind of thing it is.
\\

\addlinespace

Form (\textit{eidos}/\textit{morphē})
&
The principle through which matter is determined as
a substance of a particular kind and through which
its characteristic organization is explained.
&
Form must not be reduced to an externally imposed
shape or treated as an independently established
equivalent of identity.
\\

\addlinespace

Composite
&
The concrete substance considered as constituted
through the relation of matter and form.
&
The composite is not merely an aggregate of
independently complete constituents; its unity
requires an account of substantial organization.
\\

\addlinespace

Potentiality and actuality
&
The relation through which matter and form may be
analysed in terms of capacity and realization.
&
This relation has several applications in Aristotle
and must not be reduced to one temporal sequence
or one kind of change.
\\

\bottomrule
\end{tabularx}

\end{table}

% =============================================================================
% END OF TABLE TL01
% =============================================================================